\chapter{Datamodels}
\label{toc:datamodels}

\section{Syntax}
\label{toc:syntax}

A datamodel is defined with the clause block \texttt{datamodel}. It can
contain \texttt{pmodel}, which are ghost model fields common to all
constructors, \texttt{pinvariant}, which are global invariants for the
datamodel, and none or several constructors, defined with \texttt{pcase}.

A constructor Cons, defined with a clause block \texttt{pcase} must contain
a clause \texttt{pwhen} --- if the constructor is the last constructor defined
in the datamodel, the clause \texttt{pwhen} may be omitted --- that declares the
guard of the constructor. It can also
contain \texttt{pframe}, which are distinct local frames or self-framing local
locations, and \texttt{pinvariant}.

In function contracts, the instances of data models consumed and produced
are specified using the \texttt{produces} and \texttt{consumes} clauses.
Each global model field must be associated with a corresponding variable
or value. The instances of data models manipulated in a loop are specified
similarly using the \texttt{loop frame} clause.

The heap and the frame can be manipulated inside a function with several
clauses. The clause \texttt{heap} allows redefining the heap and the clause
\texttt{frame} allows renaming some instances of data models in the frame.
The clause \texttt{produce} is used to create a new instance of a data model in
the frame with a given constructor ; all the model fields must be associated
with a value and all the local frames must be associated with an instance of
a data model of the frame.
Conversely, the clause \texttt{consume} is used to remove an instance of a
data model of the frame with a given constructor. Similarly, to
\texttt{produce}, all the model fields and the local frame must be associated
with a value and a new instance of a data model.

\begin{figure}[htbp]
  \begin{center}
    \begin{tabular}{|ccl@{\hspace{.5cm}}l|}
    % LTeX: enabled=false
    \hline
    \rule{0pt}{2em}
    \textit{clause} & ::= & \ldots & \\
     & $\mid$ &
      \texttt{datamodel} Type \{\textit{field}, \dots, \textit{field}\} ;
    & Datamodel specification\\
    & $\mid$ &
      \texttt{consume} Name: \{ Type::Cons \textbackslash with&\\
    &&\quad.Model = Value, \dots, .Frame = Object\} ;
    & \\
    & $\mid$ &
      \texttt{produce} Name: \{ Type::Cons \textbackslash with &\\
    &&\quad.Model = Value, \dots, .Frame = Object\} ;
    & \\
    & $\mid$ &
     \texttt{consumes} CName: \{ Type \textbackslash with
    & Function contract\\
    &&\quad.Model = Value, \dots, .Model = Value\} ;&\\
    & $\mid$ &
      \texttt{produces} PName: \{ Type \textbackslash with
    & Function contract \\
    &&\quad.Model = Value,\dots, .Model = value\} ;&\\
    & $\mid$ &
      \texttt{loop frame} PName: \{ Type \textbackslash with
    & \\
    &&\quad.Model = Value,\dots, .Model = value\} ;&\\
    & $\mid$ &
      \texttt{call} \{ Fun \textbackslash with
    & Function call\\
    &&\quad.CName = - -Object, \dots, .PName = ++Object\} ;&\\
    & $\mid$ &
      \texttt{heap} &\\
    & $\mid$ &
      \texttt{frame} &\\

    \rule{0pt}{2em}
      \textit{field} & ::= &&\\
    & $\mid$ &
      \texttt{pmodel} (type) Model ; &\\
    & $\mid$ &
      \texttt{pframe} Frame: \{ Type \textbackslash with&\\
     &&\quad.Model = Value,\dots, .Frame = Object\} ; &\\
    & $\mid$ &
      \texttt{pinvariant} $\varphi$ ;&\\
    & $\mid$ &
      \texttt{pcase} Cons \{ \textit{casefield}, \dots, \textit{casefield}\} ;
    &\\

    \rule{0pt}{2em}
      \textit{casefield} & ::= &&\\
    & $\mid$ &
      \texttt{pwhen} $\varphi$ ; &\\
    & $\mid$ &
      \texttt{pframe} Frame: \{ Type \textbackslash with&\\
     &&\quad.Model = Value,\dots, .Frame = Object\} ; &\\
    & $\mid$ &
      \texttt{pinvariant} $\varphi$ ;&\\

    &&& \\
    \hline
  % LTeX: enabled=true
  \end{tabular}
  \end{center}
  \caption{Syntax of Datamodels}
  \label{fig:datamodels}
\end{figure}
